\documentclass[11pt]{article}
\usepackage[utf8]{inputenc}
\usepackage[T1]{fontenc}
\usepackage{booktabs}
\usepackage{amsmath}
\usepackage{geometry}
\geometry{margin=2.5cm}

\title{Causal Fairness Analysis Report}
\author{Generated by FairMind and an LLM}
\date{2026-08-22}

\begin{document}
\maketitle

\section{Analysis setup}

\begin{tabular}{ll}
\toprule
Dataset & diabetes\_130 \\
Protected attribute ($X$) & S1\_gender \\
Comparison & Female $\to$ Male \\
Outcome ($Y$) & T\_readmitted (YES) \\
Mediator ($W$) & time\_in\_hospital \\
Confounder ($Z$) & age \\
\bottomrule
\end{tabular}

\section{Causal effects (FairMind, exact values)}

The five effects below are computed by FairMind through exact inference on the
fitted Bayesian network. These values are final and must not be recomputed or
altered.

\begin{tabular}{lr}
\toprule
Effect & Value \\
\midrule
Total Variation (TV) & -0.018000 \\
Total Effect (TE) & -0.017200 \\
Spurious Effect (SE) & -0.000800 \\
Direct Effect (DE) & -0.000800 \\
Indirect Effect (IE, decomposition form: TE = DE + IE) & -0.016400 \\
\bottomrule
\end{tabular}

\section{Qualitative interpretation}

\subsection{Total effect and spurious component (TV, TE, SE)}

The total disparity between the two groups is relatively small, with a total variation (TV) of -0.018000. After removing the confounding effect of age (Z), the total effect (TE) is -0.017200, indicating that a small portion of the observed disparity is genuine. The spurious effect (SE) is -0.000800, which is driven by the confounder age rather than by the protected attribute itself, suggesting that the majority of the observed disparity is not due to direct discrimination.

\subsection{Direct effect (DE)}

The direct effect (DE) is -0.000800, indicating a negligible direct discrimination against the protected group. This effect is very small and does not suggest significant direct discrimination.

\subsection{Indirect effect (IE)}

The indirect effect (IE) is -0.016400, which is negative and larger in magnitude than the direct effect (DE). This indicates that the mediator channel (time in hospital) amplifies the total effect, making the overall impact more negative than the direct effect alone.

\section{Recap Questions}

\begin{enumerate}
\item Is there direct discrimination against the protected group? \textbf{LLM:} NO \quad \textbf{Expected:} NO \quad (correct)
\item Does the indirect effect through the mediator mitigate the total effect? \textbf{LLM:} NO \quad \textbf{Expected:} NO \quad (correct)
\item Does the observed total effect exceed the practical relevance threshold? \textbf{LLM:} NO \quad \textbf{Expected:} NO \quad (correct)
\item Does the spurious component (SE) contribute substantially to the observed difference? \textbf{LLM:} NO \quad \textbf{Expected:} NO \quad (correct)
\item Is the direct effect the dominant component compared to the indirect effect? \textbf{LLM:} NO \quad \textbf{Expected:} YES \quad (wrong)
\end{enumerate}

\vspace{1em}
\noindent\footnotesize
The recap answers follow deterministic threshold rules applied to the values above: direct discrimination when $|\mathrm{DE}| > 0.05$; a mediator channel counted as negligible when $|\mathrm{IE}| < 0.005$; practical relevance when $|\mathrm{TV}| > 0.05$; a substantial spurious component when $|\mathrm{SE}| / |\mathrm{TV}| > 0.1$.


\vspace{1em}
\noindent\textbf{Answer consistency:} 4/5 (80.0\%). The expected answers are derived deterministically from the FairMind values alone, through threshold rules, and were added to the document after it was generated: the model never saw them.

\end{document}